\documentclass[12pt]{article}

\usepackage[a4paper, total={6in, 8in}]{geometry}
\usepackage{textcomp}

\begin{document}
\setlength{\parindent}{0pt}

\def \t {\textrightarrow}
\def \v {\vspace{1em}}
\def \bi {\begin{itemize}}
\def \ei {\end{itemize}}
\def \s[#1] {\section*{#1}}
\def \ss[#1] {\subsection*{#1}}
\def \sss[#1] {\subsubsection*{#1}}

\s[Palazzeschi]
\ss[Chi sono?]
Finisce con i salti in banco \t delle figure ai margini con aspetto ludico, che pero nascondono sofferenza \\
Tono colloquiale \t domanda, e poi "no certo" \\
"La penna dell'anima mia" \t la parola è creazione dell'animo \t in cui è presente la follia, e il suo essere al di fuori di ogni controllo \\
La follia permette di essere se stesso \t ne abbiamo parlato con Alfieri, il folle non puo essere perseguito dal tiranno \t e anche in Alda Merini \t è un trend del 900 \\
Gioca con una colloquialita verbalizzata, che riproduce aspetto ludico ma anche aspetto della creatività: il poeta crea con la parola, il pittore con la il pennello \\
Parole chiave: follia, malinconia, nostalgia \t musico, concinnitas in "nella tastiera dell'anima mia" \t concinnitas tra strumento e arte (poesia parola, pittore tavolozza, musico tastiera) \\
La nostalgia di non poter esprimere un autenticita \t da cui diparte il rapporto dell'artista e il potere \t nel 900 la voce del poeta è difficilmente udibile \\
Lui è il giocoliere della sua anima \t il suo cuore che si rivela agli altri \t aspetto ludico del salti in banco richiama alla riflessione \\
Il testo si configura dal volto compesitivo con spazi bianchi con armonia nuova con la parola scritta, e immagine quasi futurista (che pero non è il futursimo che distrugge, qua costruisce una consapevolezza) \\
L'essere il giocoliere della sua anima, e di essere connotato con quella leggerezza che nasconde un dolore \t come la leggerezza apparentemente \\
Cavalli bianchi, lanterna e incendiario

\ss[Rio bo]
È un paese o fiume fittizio lontano dal progresso \\
È un paese microscopico, un microscosmo che crea un limite \t e che protegge da un mondo che è in progresso e crea sofferenza \\
L'autore nasce a firenze e muore a roma \t quindi sperimenta la società brulicante, e un mondo in evoluzione \\
Questo piccolo paese ha sopra una luce \t evoca il lanternino che Pirandello descrive \t semplicità apparente, ma si cela dietro la complessita \\
Il microcosmo è una protezione \t e ha una stella innamorata, che cela la domanda \t la stella guarda una grande città? \t qui invece tutti sono protetti dalla luce della stella \\
Questo è un mondo incantato \t natura che accoglie, il cipresso accoglie, e la natura è un mondo puro che non corrompe l'uomo \t filosofia di Rousseau

\v

Lo sperimentalismo di questi artisti si legge con il "Manifesto del controdolore" \t in risposta al saggio sull uomorismo di Pirandello \\
Nell'eta giovanile si piange, in quella adulta il pianto va bandito \t l'eta adulta nulla deve stupire \t il riso qua si sostituisce al pianto \t bilanciamento tra umorismo \\
Pianto giovanile si stupisce di cio che accade, che si sostituisce al riso che cela la sofferenza della vita \t il dolore è parte della vita, il controdolore è l'atteggiamento verso la vita \\
Qui ha una precisa coordinata, quella dell'eta matura \t si coglie quindi la differenza tra leopardi, che parla dell illusione e del discovrir del vero \t l'antidoto è alla fine lo stesso, il dolore è una parbola costante

\ss[Corrado Govoni]
Corrado Govoni parla della figura del palombaro poeta, senza però la profondità di ungaretti \t pubblica opere che si inseriscono nell ambito crepuscolare \\
Lascia una lirica visiva, che procedere per immagini \t parla di un futurismo agreste secondo Ermengardo \\
E rientra nel filone della quotidianetà, con Manzoni (le genti meccaniche e gli umili), nella scapigliatura (che privilegia immagine della milano del tempo), Verga (quotidiano del mondo siculo), nei crepuscolari (figure femminili non sono le dame d'annunziane)

\v

Qui invece il quotidiano è letto attraverso questo testo (Rio Bo) \\
\textbf{Elio Gioanola} \t che parla di Govoni

\ss[Il giardino - Govoni]
Un giardino primaverile \t evocazione della primavera \t ma anche autori stranieri decadenti e natura tralasciata, mondo decadente europeo \\
Borgese li definisce i crepuscolari, ma la matrice è ancora decadente \\
Evocazione di Pascoli al verso 8 (fulmine) \t abbiamo esempi che distorcono la tradizione \\
Versi centrali sono i 16, 17, 18 \t si evoca il "treno in corsa verso una citta grigia dalle strade straccioni ..." \t il mondo moderno viene letto cosi, il progresso non porta al benessere per l uomo \\
È un mondo grigio \\
L'immagine della brezza è lontana dalla tempesta e i fulmini della triade di Pascoli \t si ha un abbassamento tonale, e un linguaggio che non è mancanza di ricercatezza, ma un uso voluto dell'autore per mostrare mondo putrefatto e consunto \\
La realta non è visibile tutta \t nessuna delle certezze \\
Versi 35, 36 \t il giardino assume un volto cinereo (piogga del pineto) e funebre \t pipistrello rasenta va e viene = chiasmo, climax ascendente \\
Morte invisible che passa \t nel giardino c'è traccia di morte \t si collega al giardino del dolore \t i fiori ipnotizzano e rivelano dolore \\
La notte anche su questo giardino fa piovere una pioggia di stelle (x agosto)

\v

"Codice di Perle" \t futurista \\
"Sorelle materassi" \t vicino a mondo cinematografico \\
"Il controdolore" \\
Palazzeschi scrive anche pero un anti pioggia del pineto \t con rimbalzo andiamo, ti fermi, vai \t la costante è figure retoriche: in pascoli sono analogia e onomatopea \\
Qua invece l'onomatopea che evoca senza parole \t la fontana malata è onomatopea unica, è un unico singulto \t distorce il senso, anche la fontana fonte di acqua è malata = non senso, paroliberismo \\
\textbf{Paroliberismo} = parole in liberta, giocolieri della parola \t per Palazzeschi il \textbf{funanbolismo} della parola \t anche sperimentalismo
\end{document}
